% 6-Apr-2019
% Modified source from
% <https://www.dcl.hpi.uni-potsdam.de/media/theses/>
%
\documentclass[parskip=half, fleqn]{scrbook} % <= Druckversion: "scrbook", Bildschirmversion: "scrreprt"
% meine Einstellung
\usepackage{content/praeambel-print}

% ABOUT
\newcommand{\typ}{Projekt}
\newcommand{\autor}{Jan Unger}
\newcommand{\titel}{notizenDummyUbuntu-v02} % wird per Powershell Script ersetzt
\newcommand{\untertitel}{Mitschrift}
\newcommand{\ort}{Wuppertal}
\newcommand{\datum}{\today}
\newcommand{\website}{https://www.ju1.eu}

% DOCUMENT
%\KOMAoption{draft}{true} % <= z.B. zum "Debuggen" der Overfull-Boxes
\bibliography{content/literatur}% meine Literatur
\bibliography{content/literatur-laufen}% meine Literatur

\begin{document}
% Einband
%\pagenumbering{alph}
%\ifisbook% ju -  25-Apr-2018  - coverpage.tex 
% Modified source from
% <https://www.dcl.hpi.uni-potsdam.de/media/theses/>
%
\begin{titlepage}
	\setlength{\evensidemargin}{0.5\evensidemargin+0.5\oddsidemargin}
	\setlength{\oddsidemargin}{\evensidemargin}

	\centering

	\raisebox{-0.5\height}{\includegraphics[width=0.95\textwidth]{content/titelbild-black.pdf}}

	\vspace*{4\baselineskip}
	{\usekomafont{title}\titel\par}
	\vspace*{\baselineskip}
	{\usekomafont{author}Autor: \autor}\par

	\vfill

\end{titlepage}
\fi
%\ifisbook\cleardoubleemptypage\fi

% (Haupt-)Titelseite, Zusammenfassung, ggf. Danksagung & Inhaltsverzeichnis
\pagenumbering{roman}
% ju - 25-Apr-18 - titelpage.tex 
% Modified source from
% <https://www.dcl.hpi.uni-potsdam.de/media/theses/>
%
\begin{titlepage}
	\setlength{\evensidemargin}{0.8\evensidemargin+0.5\oddsidemargin}
	\setlength{\oddsidemargin}{\evensidemargin}

	\centering

	\raisebox{-0.5\height}{\includegraphics[width=8cm]{img/titelbild.pdf}}
	%\raisebox{-0.5\height}{\includegraphics[width=8cm]{img/titelbild.pdf}}
	%\hspace*{.2\textwidth}
	%\raisebox{-0.5\height}{\includegraphics[width=3cm]{img/logo.pdf}}
	\hfill
	\vspace{3mm}
	{\color{meinblue} \rule{\textwidth}{3pt}}


	\vspace*{4\baselineskip}
	%{\usekomafont{subject}\typ}\par
	%\vspace*{\baselineskip}
	{\usekomafont{title}\titel\par}
	\vspace*{\baselineskip}
	{\usekomafont{author}Autor: \autor}\par
	%\vspace*{\baselineskip}
	%{\usekomafont{date}\ort}\par


	\begin{flushleft}
		\vfill
		{\usekomafont{date}Version: \datum}\par
		\vspace*{\baselineskip}
		{\qrcode[hyperlink,height=1.5cm]{\website}}\par
	\end{flushleft}

	\setcounter{page}{1}

\end{titlepage}

%\ifisbook\cleardoubleemptypage\fi% ju - 25-Apr-18 - zusammenfassung.tex 
% Modified source from
% <https://www.dcl.hpi.uni-potsdam.de/media/theses/>
%
% deutsche Zusammenfassung
\null\vfil
\begin{otherlanguage}{ngerman}
\begin{center}\textsf{\textbf{\abstractname}}\end{center}

\noindent Dies hier ist ein Blindtext zum Testen von Textausgaben. Wer diesen Text liest,
ist selbst schuld. Der Text gibt lediglich den Grauwert der Schrift an.
Ist das wirklich so? Ist es gleichgültig, ob ich schreibe: "Dies ist ein Blindtext"
oder "Huardest gefburn"? Kjift - mitnichten! Ein Blindtext bietet mir wichtige Informationen.

\end{otherlanguage}
\vfil\null
% example
%\ifisbook\cleardoubleemptypage\fi% ju - 25-Apr-18  - danksagung.tex 
% Modified source from
% <https://www.dcl.hpi.uni-potsdam.de/media/theses/>
%
\vspace*{\fill}
\begin{center}\textsf{\textbf{Danksagung}}\end{center}

\noindent Dies hier ist ein Blindtext zum Testen von Textausgaben. Wer diesen Text liest,
ist selbst schuld. Der Text gibt lediglich den Grauwert der Schrift an.
Ist das wirklich so? Ist es gleichgültig, ob ich schreibe: "Dies ist ein Blindtext"
oder "Huardest gefburn"? Kjift - mitnichten! Ein Blindtext bietet mir wichtige Informationen.

\vspace*{\fill}
% example
\tableofcontents
\cleardoublepage

% Textteil
\pagenumbering{arabic}

% Inhalt
% --------------------------------------------------
%
% ---------------------------------
\chapter{neu}
%----------
% \section{ }
% \subsection{ }\label{ }\index{ }
%----------

\hypertarget{test}{%
\section{Test}\label{test}}

\hypertarget{quellcode}{%
\subsection{Quellcode}\label{quellcode}}

(\autoref{code:Codename } Codename). % Codeverweis =

% Quellcode
\lstset{language=Bash} % C, [LaTeX]TeX, Bash, Python
\begin{lstlisting}[numbers=left, frame=l, framerule=0.1pt,%
% -----------
	caption={Codename }, % Caption
	label={code:Codename }  % Label
]% ---------=

  # Überschrift
  ## Überschrift 2
  ### Überschrift 3
\end{lstlisting}

\hypertarget{bild}{%
\subsection{Bild}\label{bild}}

% Bild Referenz
(\autoref{pic:logo} logo.pdf ). % Bildverweis = logo.pdf

\begin{figure}[!hp]% hier
	\centering
	\includegraphics[width=0.8\textwidth]{img/logo.pdf}
	% ---------
	\caption[logo.pdf ]{logo.pdf } % Caption
	\label{pic:logo}	% Referenz
	% ---------
\end{figure}

\hypertarget{tabelle}{%
\subsection{Tabelle}\label{tabelle}}

(\autoref{tab:Tabellenname } Tabellenname ). % Tabellenverweis = Tabellenname

% Tabelle
\begin{table}[!hp] % hier
	\centering
	%\setlength{\tabcolsep}{5mm} % Spaltenlänge fest
	\rowcolors{1}{}{lightgray!20}  % Farbe
	%\begin{tabularx}{\textwidth}{XX} % auto. Spaltenumbruch
	\begin{tabular} []{rll}
	\toprule % Inhalt
	% ---------
\textbf{Nr.} & \textbf{Begriffe} & \textbf{Erklärung}\\
	\midrule
1 & a1 & a2\\
2 & b1 & b2\\
3 & c1 & c2\\
	\bottomrule
	%\end{tabularx}
	\end{tabular}
	\caption[Tabellenname ]{Tabellenname }	% Caption
	\label{tab:Tabellenname }	% Referenz
	%-----------
\end{table}
%-----------



 % wird per Powershell Script ersetzt
%
% --------------------------------------------------
%\chapter{Kapitel}
%\chapter{Readme}
%\include{content/Readme}
%\chapter{MarkdownSpickzettel}
%%============================
% \section{ }
% \subsection{ }\label{ }\index{ }
%============================

\section{Markdown - Spickzettel}\label{markdown-syntax}

\subsection{Überschrift}\label{ueberschrift}

% Quellcode
\lstset{language=Bash} % C, [LaTeX]TeX, Bash, Python
\begin{lstlisting}[numbers=left, frame=l, framerule=0.1pt,%
% ======================
	caption={ }, % Caption
	label={code: }  % Label
]% =====================

    # Überschrift
    ## Überschrift 2
    ### Überschrift 3
\end{lstlisting}

\subsection{Bild}\label{bild}

% Bild Referenz
(\autoref{pic:logo1} logo.pdf). % Bildverweis = logo.pdf

\begin{figure}[!hp]% hier
	\centering
	\includegraphics[width=0.8\textwidth]{img/logo.pdf}
	% ====================
	\caption[logo.pdf ]{logo.pdf } % Caption
	\label{pic:logo1}	% Referenz
	% ====================
\end{figure}

% Quellcode
\lstset{language=Bash} % C, [LaTeX]TeX, Bash, Python
\begin{lstlisting}[numbers=left, frame=l, framerule=0.1pt,%
% ======================
	caption={ }, % Caption
	label={code: }  % Label
]% =====================

    # bild
    ![Logo](img/logo.pdf)
\end{lstlisting}

\subsection{Tabelle}\label{tabelle}

(\autoref{tab:Tabelle } ). % Tabellenverweis = Tabelle

% Tabelle
\begin{table}[!hp] % hier
	\centering
	%\setlength{\tabcolsep}{5mm} % Spaltenlänge fest
	\rowcolors{1}{}{lightgray!20}  % Farbe
	%\begin{tabularx}{\textwidth}{XX} % auto. Spaltenumbruch
	\begin{tabular} []{rll}
	\toprule % Inhalt
	% ---------
\textbf{Nr.} & \textbf{Begriffe} & \textbf{Erklärung}\\
	\midrule
1 & a1 & a2\\
2 & b1 & b2\\
3 & c1 & c2\\
4 & a1 & a2\\
	\bottomrule
	%\end{tabularx}
	\end{tabular}
	\caption[Tabelle ]{Tabelle }	% Caption
	\label{tab:Tabelle }	% Referenz
	%============================
\end{table}

% Quellcode
\lstset{language=Bash} % C, [LaTeX]TeX, Bash, Python
\begin{lstlisting}[numbers=left, frame=l, framerule=0.1pt,%
% ======================
	caption={ }, % Caption
	label={code: }  % Label
]% =====================

    # tabelle
    |**Nr.**|**Begriffe**|**Erklärung**|
    |--:|:-----|:----|
    | 1     | a1         | a2                |
    | 2     | b1         | b2              |
    | 3     | c1         | c2              |
    | 4     | a1         | a2              |
\end{lstlisting}

\subsection{Mathe}\label{mathe}

$[ V ] = [ \Omega ] \cdot [ A ]$ o. $U = R \cdot I$ o.
$R = \frac{U}{I}$

% Quellcode
\lstset{language=Bash} % C, [LaTeX]TeX, Bash, Python
\begin{lstlisting}[numbers=left, frame=l, framerule=0.1pt,%
% ======================
	caption={ }, % Caption
	label={code: }  % Label
]% =====================

    # Mathe
    $[ V ] = [ \Omega ] \cdot [ A ]$ o. $U = R \cdot I$ o. $R = \frac{U}{I}$
\end{lstlisting}

\textbf{Matheumgebung:}

\begin{align*}
    \sum_{i=1}^5 a_i = a_1 + a_2 + a_3 + a_4 + a_5
\end{align*}

% Quellcode
\lstset{language=Bash} % C, [LaTeX]TeX, Bash, Python
\begin{lstlisting}[numbers=left, frame=l, framerule=0.1pt,%
% ======================
	caption={ }, % Caption
	label={code: }  % Label
]% =====================

    # Matheumgebung
    \begin{align*}
        \sum_{i=1}^5 a_i = a_1 + a_2 + a_3 + a_4 + a_5
    \end{align*}
\end{lstlisting}

\subsection{Absätze}\label{absaetze}

Dies hier ist ein Blindtext zum Testen von Textausgaben. Wer diesen Text
liest, ist selbst schuld. Der Text gibt lediglich den Grauwert der
Schrift an. Ist das wirklich so? Ist es gleichgültig, ob ich schreibe:
\frqq Dies ist ein Blindtext\flqq\, oder \frqq Huardest gefburn\flqq\,? Kjift -
mitnichten! Ein Blindtext bietet mir wichtige Informationen.

Dies hier ist ein Blindtext zum Testen von Textausgaben. Wer diesen Text
liest, ist selbst schuld. Der Text gibt lediglich den Grauwert der
Schrift an. Ist das wirklich so? Ist es gleichgültig, ob ich schreibe:
\frqq Dies ist ein Blindtext\flqq\, oder \frqq Huardest gefburn\flqq\,? Kjift -
mitnichten! Ein Blindtext bietet mir wichtige Informationen.

\subsection{Listen}\label{listen}

\textbf{ungeordnete Liste}

\begin{itemize}% Liste Punkt
\item
  a
\item
  b

  \begin{itemize}% Liste Punkt
  \item
    bb
  \end{itemize}
\item
  c
\end{itemize}

% Quellcode
\lstset{language=Bash} % C, [LaTeX]TeX, Bash, Python
\begin{lstlisting}[numbers=left, frame=l, framerule=0.1pt,%
% ======================
	caption={ }, % Caption
	label={code: }  % Label
]% =====================

    # ungeordnete Liste
    - a
    - b
        - bb
    - c
\end{lstlisting}

\textbf{Sortierte Liste}

\begin{enumerate}% Liste 1) oder a)
\item
  eins
\item
  zwei
\item
  drei
\end{enumerate}

% Quellcode
\lstset{language=Bash} % C, [LaTeX]TeX, Bash, Python
\begin{lstlisting}[numbers=left, frame=l, framerule=0.1pt,%
% ======================
	caption={ }, % Caption
	label={code: }  % Label
]% =====================

    # Sortierte Liste
    1. eins
    2. zwei
    3. drei
\end{lstlisting}

\textbf{Sortierte Liste}

\begin{enumerate}% Liste 1) oder a)
\def\labelenumi{\alph{enumi})}
\item
  a
\item
  b
\item
  c
\end{enumerate}

(\autoref{code:Liste } Liste). % Codeverweis = Liste

% Quellcode
\lstset{language=Bash} % C, [LaTeX]TeX, Bash, Python
\begin{lstlisting}[numbers=left, frame=l, framerule=0.1pt,%
% ======================
	caption={Liste }, % Caption
	label={code:Liste }  % Label
]% =====================

    # Sortierte Liste
    a) a
    b) b
    c) c
\end{lstlisting}

\subsection{Anführungszeichen}\label{anfuehrungszeichen}

\frqq Anführungszeichen\flqq\, oder \flqq Anführungszeichen\frqq   oder
\frqq Anführungszeichen\flqq

% Quellcode
\lstset{language=Bash} % C, [LaTeX]TeX, Bash, Python
\begin{lstlisting}[numbers=left, frame=l, framerule=0.1pt,%
% ======================
	caption={ }, % Caption
	label={code: }  % Label
]% =====================

    # Anführungszeichen
    "Anführungszeichen" oder
    \flqq Anführungszeichen\frqq   oder
    \frqq Anführungszeichen\flqq
\end{lstlisting}

\subsection{Quellenangabe}\label{quellenangabe}

Quelle \textcite{schlosser_latex:2016}

Quelle Text\footfullcite{schlosser_latex:2016}

Quelle \cite{schlosser_latex:2016}

% Quellcode
\lstset{language=Bash} % C, [LaTeX]TeX, Bash, Python
\begin{lstlisting}[numbers=left, frame=l, framerule=0.1pt,%
% ======================
	caption={ }, % Caption
	label={code: }  % Label
]% =====================

    # Quellenangabe
    Quelle \textcite{schlosser_latex:2016}
    Quelle Text\footfullcite{schlosser_latex:2016}
    Quelle \cite{schlosser_latex:2016}
\end{lstlisting}
%============================
%============================
%============================

%\chapter{LatexSpickzettel}
%% ju - 21-Nov-18 - kap1.tex 
%\chapter{Kapitel}
%
\section{\LaTeX - Spickzettel}\label{sec:LaTeX-Spickzettel}

\subsection{Blindtext}\label{sec:blindtext}

\blindtext[1] \\ \blindtext[1]

\subsection{Flattersatz versus Blocksatz}\label{sec:Flattersatz-versus-Blocksatz}

\begin{flushleft}
\blindtext[1]
\end{flushleft}

\begin{flushright}
\blindtext[1]
\end{flushright}

\begin{center}
\blindtext[1]
\end{center}

\newpage % ++++++++++++++

\subsection{Gliederung}

% Quellcode Referenz
(\autoref{code:gliederung} Gliederung).
% Quellcode
\lstset{language=[LaTeX]TeX} % C, [LaTeX]TeX, Bash, Python
\begin{lstlisting}[numbers=left, frame=l, framerule=0.1pt,%
% ======================
  caption={Gliederung},    % Caption anpassen!
  label={code:gliederung}  % Label anpassen!
]% =====================

	% artikel:
	\subsection \subsubsection \paragraph \subparagraph

	% report:
	\chapter \subsection \subsubsection \paragraph \subparagraph

	% * keine Nummerierung
	\subsection*{eins}

	% Sprungmarke
	\label{sec:eins}
\end{lstlisting}

\newpage % ++++++++++++++

\subsection{Quellcode}

% Quellcode Referenz
(\autoref{code:dummyCode} dummyCode).  % Anpassen!
% Quellcode
\lstset{language=[LaTeX]TeX} % C, [LaTeX]TeX, Bash, Python
\begin{lstlisting}[numbers=left, frame=l, framerule=0.1pt,%
% ======================
  caption={dummyCode},    % Caption anpassen!
  label={code:dummyCode}  % Label anpassen!
]% ======================

	/* Quellcode */
	\verb|Code|
\end{lstlisting}

% Quellcode Ausgabe Referenz
(\autoref{code:dummyCode-out} dummyCode Ausgabe).  % Anpassen!
% Quellcode
\lstset{language=[LaTeX]TeX} % C, [LaTeX]TeX, Bash, Python
\begin{lstlisting}[numbers=left, frame=l, framerule=0.1pt,%
% ======================
  caption={dummyCode Ausgabe},    % Caption anpassen!
  label={code:dummyCode-out}  % Label anpassen!
]% ======================

	% Quellcode Referenz
	(\autoref{code:dummyCode} dummyCode).  % Anpassen!
	% Quellcode
	\lstset{language=[LaTeX]TeX} % C, [LaTeX]TeX, Bash, Python
	\begin{lstlisting}[numbers=left, frame=l, framerule=0.1pt,%
	% ======================
		caption={dummyCode},    % Caption anpassen!
		label={code:dummyCode}  % Label anpassen!
	]% ======================

		/* Quellcode */
	%end{lstlisting}
\end{lstlisting}

% Quellcode Referenz
(\autoref{code:hallo-ext} hallo.c).    % codename =
% Quellcode aus ext. Datei
	\lstset{language=C}% C++, [LaTeX]TeX, Bash, Python
	\lstinputlisting[numbers=left, frame=l, framerule=0.1pt,
	% =====================
	caption={Quellcode in C, hallo.c},    % Caption  anpassen!
	label={code:hallo-ext}  % Referenz anpassen!
	% =====================
	]{content/hallo.c}% ext. Datei

\newpage % ++++++++++++++

\subsection{Querverweise-Referenzen}\label{sec:quer-ref}

(\autoref{sec:schriftstile} Schriftstile).

(\autoref{sec:listen} Listen).

(\autoref{pic:tux1} Tux1).

\begin{enumerate}
	\item zuerst
	\item \label{item:folge} folgend
	\item abschließend
\end{enumerate}

Wir verweisen auf ein Listenelement \autoref{item:folge}.

% Quellcode Ausgabe Referenz
(\autoref{code:quer-ref-out} Querverweise-Referenzen).% Anpassen!
% Quellcode
\lstset{language=[LaTeX]TeX} % C, [LaTeX]TeX, Bash, Python
\begin{lstlisting}[numbers=left, frame=l, framerule=0.1pt,%
% ======================
	caption={Quellcode in LaTeX, Querverweise-Referenzen},% Anpassen!
	label={code:quer-ref-out}%
]% -----------------------

	(\autoref{sec:schriftstile} Schriftstile).
	(\autoref{sec:listen} Listen).
	(\autoref{pic:tux1} Tux1).
	\begin{enumerate}
		\item zuerst
		\item \label{item:folge} folgend
		\item abschließend
	\end{enumerate}
	Wir verweisen auf ein Listenelement \autoref{item:folge}.
\end{lstlisting}

% Tabelle Referenz
(\autoref{tab:label-verweis} Label Querverweis).
% Tabelle
\begin{table}[!hb] % hier einfügen
	\centering
	%\setlength{\tabcolsep}{5mm}      % Spaltenlänge fest
  \rowcolors{1}{}{lightgray!20}     % Farbe
	%\begin{tabularx}{\textwidth}{XX} % auto. Spaltenumbruch
	\begin{tabular} {ll}
		\toprule % Inhalt
		% ---------------------
			Abk. & Beschreibung \\
		\midrule
			sec  & für alle Gliederungsebenen \\
			cha  & oder chap für Kapitel (es kann aber auch sec verwendet werden) \\
			part & für Teile eines Buches (ebenso sec möglich) \\
			fig  & für Abbildungen \\
			tab  & für Tabellen \\
			item & für Aufzählungspunkte \\
			eqn  & für Gleichungen  \\
			fn   & für Fußnoten \\
			code & Listing \\
			pic  & Grafik \\
		% ----------------------
		\bottomrule
	\end{tabular}
	% =======================
		\caption{Label Querverweis }	  % Caption anpassen!
		\label{tab:label-verweis}	      % Referenz anpassen!
	% =======================
\end{table}

\newpage % ++++++++++++++

\subsection{Zitieren}\label{sec:zitieren}

\begin{quote}
Ein schönes Zitat von einem schlauen Menschen steht den meisten Dokumenten
gut zu Gesicht.
\end{quote}

Fussnote\footnote{Fussnote}.

Google\footnote{\url{https://www.google.de/}}

Anfang ">Anführungszeichen."< Ende

Anfang "`Anführungszeichen."' Ende

Anfang \flqq Anführungszeichen.\frqq Ende

\LaTeX Quelle \textcite{schlosser_latex:2016}.

% Quellcode Referenz
(\autoref{code:zitieren-out} Zitieren).
% Quellcode
\lstset{language=[LaTeX]TeX} % C, [LaTeX]TeX, Bash, Python
\begin{lstlisting}[numbers=left, frame=l, framerule=0.1pt,%
% ======================
	caption={Quellcode in LaTeX, Zitieren},% Anpassen!
	label={code:zitieren-out}%
]% -----------------------

	\begin{quote}
		Ein schönes Zitat von einem schlauen Menschen steht
		den meisten Dokumenten gut zu Gesicht.
	\end{quote}
	Fussnote\footnote{Fussnote}.
	Anfang ">Anführungszeichen."< Ende
	Anfang "`Anführungszeichen."' Ende
	Anfang \flqq Anführungszeichen.\frqq Ende
	\LaTeX Quelle \textcite{schlosser_latex:2016}.
\end{lstlisting}

\newpage % ++++++++++++++

\subsection{Links}\label{sec:links}\index{Links}

Darstellung einer klickbaren URL: \url{https://www.google.de/}

Text, der auf eine Webseite linkt: \href{/}{Meine Webseite}

Emailadresse verlinken: \href{mailto:jan@ju1.eu}{Meine E-Mail-Adresse}

auf lokale Datei verlinken: \href{run:/content/hallo.c}{lokale Datei}

PDF einbinden: \includepdf[pagecommand={\thispagestyle{headings}},
	noautoscale=true,width=0.9\textwidth,offset=0cm -1cm]{content/titelbild.pdf}

PDF einbinden: \\

\includegraphics[width=0.9\textwidth]{content/titelbild.pdf}


(\autoref{code:links-out} Links).
% Quellcode
\lstset{language=[LaTeX]TeX} % C, [LaTeX]TeX, Bash, Python
\begin{lstlisting}[numbers=left, frame=l, framerule=0.1pt,%
% ======================
	caption={Quellcode in LaTeX, Links},% Anpassen!
	label={code:links-out}
]%======================

	Darstellung einer klickbaren URL: \url{https://www.google.de/}
	Text, der auf eine Webseite linkt: \href{/}{Meine Webseite}
	Emailadresse verlinken: \href{mailto:jan@ju1.eu}{Meine E-Mail-Adresse}
	auf lokale Datei verlinken: \href{run:/content/hallo.c}{lokale Datei}
	PDF einbinden: \includepdf[pagecommand={\thispagestyle{headings}},
		noautoscale=true,width=0.9\textwidth,offset=0cm -1cm]{content/titelbild.pdf}
	PDF einbinden: \\
		\includegraphics[width=0.9\textwidth]{content/titelbild.pdf}
\end{lstlisting}

\newpage % ++++++++++++++

\subsection{Farbe}

Text \textcolor{red}{rot} Text

Text \colorbox{hellesbrombeer}{hellesbrombeer} Text

\colorbox{red!10!white}{10 \% rot, Rest weiß}

\textcolor{meinorange}{farbiger Text}
\textcolor{meinblue}{farbiger Text}
\textcolor{meinred}{farbiger Text}
\textcolor{magenta}{farbiger Text}

\wichtig[meinblue]{wichtiger farbiger Text}
\wichtig[meinred]{wichtiger farbiger Text}
\wichtig[meinorange]{wichtiger farbiger Text}
\wichtig[magenta]{wichtiger farbiger Text}

% Farbbox
\colorbox{meingrey}{Text}
\colorbox{meinorange}{Text}

% bunter Rahmen um eine Formel
\fcolorbox{meinblue}{meingrey}{$a^{2} + b^{2} = c^{2}$}

% Quellcode Referenz
(\autoref{code:farbe-out} Farbe).
% Quellcode
\lstset{language=[LaTeX]TeX} % C, [LaTeX]TeX, Bash, Python
\begin{lstlisting}[numbers=left, frame=l, framerule=0.1pt,%
% ======================
	caption={Quellcode in LaTeX, Farbe},% Anpassen!
	label={code:farbe-out}
]%======================

	\textcolor{meingreen}{farbiger Text}
	\textcolor{meinblue}{farbiger Text}
	\textcolor{meinred}{farbiger Text}

	\wichtig[meinblue]{wichtiger farbiger Text}
	\wichtig[meinred]{wichtiger farbiger Text}
	\wichtig[meingreen]{wichtiger farbiger Text}

	% Farbbox
	\colorbox{meingrey}{Text}
	\colorbox{meinorange}{Text}

	% bunter Rahmen um eine Formel
	\fcolorbox{meinblue}{meingrey}{$a^{2} + b^{2} = c^{2}$}
\end{lstlisting}

\subsection{Hinweis - Box}

\begin{hinweis}\label{sec:hinweis}
  Als bekanntester Blindtext gilt der Text "Lorem ipsum", der seinen Ursprung im 16. Jahrhundert haben soll.
  Lorem ipsum ist in einer pseudo-lateinischen Sprache verfasst, die ungefähr dem "natürlichen" Latein entspricht.
  In Ihm finden sich eine Reihe realer lateinischer Wörter.
\end{hinweis}

\myInfoBox{
  Als bekanntester Blindtext gilt der Text "Lorem ipsum", der seinen Ursprung im 16. Jahrhundert haben soll.
  Lorem ipsum ist in einer pseudo-lateinischen Sprache verfasst, die ungefähr dem "natürlichen" Latein entspricht.
  In Ihm finden sich eine Reihe realer lateinischer Wörter.
}


% Quellcode Referenz
(\autoref{code:hinweis-infobox-out} Hinweis, Infobox).
% Quellcode
\lstset{language=[LaTeX]TeX} % C, [LaTeX]TeX, Bash, Python
\begin{lstlisting}[numbers=left, frame=l, framerule=0.1pt,%
% ======================
	caption={Quellcode in LaTeX, Hinweis, Infobox},% Anpassen!
	label={code:hinweis-infobox-out}
]%======================

	\begin{hinweis}
    Text
	\end{hinweis}

	\myInfoBox{
    Text
	}
  \mybox{
    Text
  }
\end{lstlisting}

\newpage % ++++++++++++++

\subsection{Tabellen}\label{sec:tabellen}

\begin{tabular}{ccc}
	\toprule
		Leistung & 45 & kWh \\
	\midrule
		Hubraum & $1234$ & $cm^3$ \\
		Preis & 23499 & Euro \\
	\bottomrule
\end{tabular}\\

Text

% Tabellen Referenz
(\autoref{tab:dummyTabelle} dummyTabelle).
% Tabelle
\begin{table}[!hb] % hier
	\centering
	%\setlength{\tabcolsep}{5mm}      % Spaltenlänge fest
  \rowcolors{1}{}{lightgray!20}     % Farbe
	%\begin{tabularx}{\textwidth}{XX} % auto. Spaltenumbruch
	\begin{tabular} {ll}
		\toprule
		% -----------------------
			\textbf{A} & \textbf{B} \\
		\midrule
			a1 & a2 \\
			b1 & b2 \\
			c1 & c2 \\
		% -----------------------
		\bottomrule
	%\end{tabularx}
	\end{tabular}
	% =======================
		\caption{dummyTabelle}          % Caption anpassen!
		\label{tab:dummyTabelle}        % Referenz anpassen!
	% =======================
\end{table}

% Quellcode Referenz
(\autoref{code:dummyTabelle-out} dummyTabelle).
% Quellcode
\lstset{language=[LaTeX]TeX} % C, [LaTeX]TeX, Bash, Python
\begin{lstlisting}[numbers=left, frame=l, framerule=0.1pt,%
% ======================
	caption={Quellcode in LaTeX, dummyTabelle},% Anpassen!
	label={code:dummyTabelle-out}%
]%======================

	% Tabellen Referenz
	(\autoref{tab:dummyTabelle} dummyTabelle).
	% Tabelle
	\begin{table}[!hb] % hier
		\centering
		%\setlength{\tabcolsep}{5mm}      % Spaltenlänge fest
		\rowcolors{1}{}{lightgray!20}     % Farbe
		%\begin{tabularx}{\textwidth}{XX} % auto. Spaltenumbruch
		\begin{tabular} {ll}
			\toprule
			% -----------------------
				\textbf{A} & \textbf{B} \\
			\midrule
				a1 & a2 \\
				b1 & b2 \\
				c1 & c2 \\
			% -----------------------
			\bottomrule
		%\end{tabularx}
		\end{tabular}
		% =======================
			\caption{dummyTabelle}          % Caption anpassen!
			\label{tab:dummyTabelle}        % Referenz anpassen!
		% =======================
	\end{table}
\end{lstlisting}

% Tabellen Referenz
(\autoref{tab:spalte-fest} Spaltenlänge fest).
% Tabelle
\begin{table}[!hb] % hier
	\centering
	\setlength{\tabcolsep}{5mm}       % Spaltenlänge fest
  \rowcolors{1}{}{lightgray!20}     % Farbe
	%\begin{tabularx}{\textwidth}{XX} % auto. Spaltenumbruch
	\begin{tabular} {ll}
		\toprule
		% -----------------------
	    \textbf{A} & \textbf{B} \\
	  \midrule
			a1 & a2 \\
			b1 & b2 \\
			c1 & c2 \\
		% -----------------------
		\bottomrule
	%\end{tabularx}
	\end{tabular}
	% -----------------------
		\caption{Spaltenlänge fest}     % Caption anpassen!
		\label{tab:spalte-fest}		      % Referenz anpassen!
	% -----------------------
\end{table}

(Tabelle Longtable).
% Longtable
\rowcolors{1}{}{lightgray!20} % Farbe
\begin{longtable}{ll}% Spaltenanzahl, l,r,c,p,X
	\toprule
	\textbf{A} & \textbf{B} \\
	\midrule
	\endfirsthead
	\toprule
	% Tab.-Fortsetzung
	\textbf{A} & \textbf{B} \\
	\midrule
	\endhead
	% Inhalt
	a1 & a2 \\
	b1 & b2 \\
	c1 & c2 \\
	a1 & a2 \\
	b1 & b2 \\
	c1 & c2 \\
	a1 & a2 \\
	b1 & b2 \\
	c1 & c2 \\
	a1 & a2 \\
	b1 & b2 \\
	c1 & c2 \\
	c1 & c2 \\
	a1 & a2 \\
	b1 & b2 \\
	c1 & c2 \\
	\bottomrule
\end{longtable}

\newpage % ++++++++++++++

\subsection{Abbildungen}\label{sec:abbildungen}

% Bild Referenz
(\autoref{pic:tux1} Linux Pinguin Tux).
% Bild
\begin{figure}[!hb] % hier
	\centering
  \includegraphics[width=0.3\textwidth]{img/logo.pdf}
	% ====================
	\caption[Tux]{">Ein wohlgenährter, glücklicher, rundlicher Pinguin, ist das
									offizielle Maskottchen des freien Betriebssystemkerns Linux."<
	\newline {(Quelle:~Wikipedia)}}% Caption  anpassen!
	\label{pic:tux1}% Referenz anpassen!
	% ===================
\end{figure}

\includegraphics[height=2cm,width=3cm]{img/logo.pdf}
\includegraphics[angle=45, scale=.2]{img/logo.pdf}

% Bild Referenz
(\autoref{pic:dummyAbb} dummyAbb).    % Anpassen!
% Bild
\begin{figure}[!hb]% hier
	\centering
  \includegraphics[width=0.3\textwidth]{img/logo.pdf}
	% ====================
	\caption{dummyAbb}% Caption  anpassen!
	\label{pic:dummyAbb}% Referenz anpassen!
	% ===================
\end{figure}

% Bild Referenz
(\autoref{pic:bildDrehen} Drehen um 45 Grad).    % Anpassen!
% Bild
\begin{figure}[!hb]% hier
	\centering
  \includegraphics[angle=45,width=0.3\textwidth]{img/logo.pdf}
	% ====================
	\caption{Drehen um 45 Grad}% Caption  anpassen!
	\label{pic:bildDrehen}% Referenz anpassen!
	% ===================
\end{figure}

\newpage % ++++++++++++++


% Quellcode Referenz
(\autoref{code:dummyAbb-out} dummyAbb).
% Quellcode
\lstset{language=[LaTeX]TeX} % C, [LaTeX]TeX, Bash, Python
\begin{lstlisting}[numbers=left, frame=l, framerule=0.1pt,%
% ======================
	caption={Quellcode in LaTeX, dummyAbb},% Anpassen!
	label={code:dummyAbb-out}%
]%======================

	% Bild Referenz
	(\autoref{pic:dummyAbb} dummyAbb).    % Anpassen!
	% Bild
	\begin{figure}[!hb]% hier
		\centering
		\includegraphics[width=0.3\textwidth]{img/logo.pdf}
		% ====================
		\caption{dummyAbb}% Caption  anpassen!
		\label{pic:dummyAbb}% Referenz anpassen!
		% ===================
	\end{figure}
\end{lstlisting}

\newpage % ++++++++++++++

\subsection{Scalieren}\index{Skalieren}

Inhalt n-fach scalieren

\vspace{10mm}

\scalebox{3}{Text}
\scalebox{4}{Text}
\scalebox{5}{Text}

\subsection{Rotieren}

Inhalt rotieren - Wert in Grad

\vspace{10mm}

\rotatebox{45}{Text}
\rotatebox{90}{Text}
\rotatebox{180}{Text}


% Quellcode Referenz
(\autoref{code:scal-rot-out} Scalieren und Rotieren).
% Quellcode
\lstset{language=[LaTeX]TeX} % C, [LaTeX]TeX, Bash, Python
\begin{lstlisting}[numbers=left, frame=l, framerule=0.1pt,%
% ======================
	caption={Quellcode in LaTeX, Scalieren und Rotieren},% Anpassen!
	label={code:scal-rot-out}%
]%======================

	% Inhalt n-fach scalieren
	\scalebox{3}{Text}
	\scalebox{4}{Text}
	\scalebox{5}{Text}

	% Inhalt rotieren - Wert in Grad
	\rotatebox{45}{Text}
	\rotatebox{90}{Text}
	\rotatebox{180}{Text}
\end{lstlisting}

\newpage % ++++++++++++++

\subsection{Gliederung in Kapitel und Abschnitte}\label{sec:Gliederung-Kapitel-Abschnitte}

(\autoref{tab:Gliederung-Kapitel-Abschnitte} Gliederung in Kapitel und Abschnitte).
% Tabelle
\begin{table}[!hb]% hier
	\centering
	%\setlength{\tabcolsep}{5mm}      % Spaltenlänge fest
  \rowcolors{1}{}{lightgray!20}     % Farbe
	%\begin{tabularx}{\textwidth}{XX} % auto. Spaltenumbruch
	\begin{tabular} {ll}
		\toprule % Inhalt
		% -----------------------
			\verb|#| & \textbf{Beschreibung} \\
	  \midrule
			\verb|\chapter{}      | & Ein Kapitel\\
			\verb|\section{}      | & Ein Abschnitt\\
			\verb|\subsection{}   | & Ein Unterabschnitt\\
			\verb|\subsubsection{}| & Ein Unter-Unterabschnitt\\
			\verb|\paragraph{}    | & Ein Absatz\\
			\verb|\subparagraph{} | & Ein Unterabsatz\\
			\verb|\subsection*{}  | & Ein unnummerierter Abschnitt\\
			\verb|\subsection[Kurzer Titel]{}| & langer Abschnittstitel\\
		% -----------------------
		\bottomrule
	%\end{tabularx}
	\end{tabular}
	% =========================
		\caption{Gliederung in Kapitel und Abschnitte}	% Caption anpassen!
		\label{tab:Gliederung-Kapitel-Abschnitte}	      % Referenz anpassen!
	% =========================
\end{table}

\subsection{Schriftstile}\label{sec:schriftstile}

\emph{kursiv}
\textrm{Antiqua}, \textsf{Grotesk}, \texttt{Maschinenschrift},
\textmd{normal}, \textbf{breiter}, \textup{aufrecht}, \textsl{geneigt},
\textit{kursiv}, \textsc{Kapitaelchen}

\subsection{Schriftgrößen}

\tiny{winzig}, \scriptsize{sehr klein}, \footnotesize{klein},
\small{klein}, , \large{gross}, \Large{groesser},
\LARGE{ganz gross}, \huge{riesig}, \Huge{gigantisch} \normalsize{normal}

\subsection{Wortabstände}\label{sec:Wortabstaende}

(\autoref{tab:wortabstand} Wortabstände).
% Tabelle
\begin{table}[!hb] % hier
	\centering
	%\setlength{\tabcolsep}{5mm}      % Spaltenlänge fest
  \rowcolors{1}{}{lightgray!20}     % Farbe
	%\begin{tabularx}{\textwidth}{XX} % auto. Spaltenumbruch
	\begin{tabular} {ll}
		\toprule % Inhalt
		% -----------------------
			\verb|#| & \textbf{Beschreibung} \\
		\midrule
			\verb|\| & erzeugt Leerstelle \\
			\verb|\@| & kennzeichnet einen Punkt als Satzende \\
			\verb|~| & erzeugt nicht umbrechbare Leerstelle \\
			\verb|\,| & erzeugt nicht umbrechbare Leerstelle \\
			\verb|\quad| & erzeugt einfach vergrößerten Abstand \\
			\verb|\qquad| & erzeugt vierfach vergrößerten Abstand \\
			\verb|\hspace{1cm}| & erzeugt Abstand von 1cm Breite \\
			\verb|\hfill| & fügt so viel Leerraum ein wie möglich \\
			\verb|\smallskip| & vertikaler Abstand\\
			\verb|\medskip| & \\
			\verb|\bigskip| & \\
			\verb|\vspace{1cm}| & \\
			\verb|\vfill| & \\
		% -----------------------
		\bottomrule
	%\end{tabularx}
	\end{tabular}
	% ======================
		\caption{Wortabstände }	% Caption anpassen!
		\label{tab:wortabstand}	% Referenz anpassen!
	% ======================
\end{table}

\subsection{Logische Textauszeichnung}\label{sec:LogischeTextauszeichnung}

(\autoref{tab:Textauszeichnung} Logische Textauszeichnung).
% Tabelle
\begin{table}[!hb] % hier
	\centering
	%\setlength{\tabcolsep}{5mm}      % Spaltenlänge fest
  \rowcolors{1}{}{lightgray!20}     % Farbe
	%\begin{tabularx}{\textwidth}{XX} % auto. Spaltenumbruch
	\begin{tabular} {ll}
		\toprule % Inhalt
		% -----------------------
			\verb|#| & \textbf{Beschreibung} \\
		\midrule
			\verb|\emph{Hervorhebung}| & \emph{Hervorhebung} \\
			\verb|\url{http://www.dante.de/}| & \url{http://www.dante.de/} \\
			\verb|\href{/}{Meine Webseite} |& \href{/}{Meine Webseite} \\
			\verb|\href{mailto:jan@ju1.eu}{jan@ju1.eu} |& \href{mailto:jan@ju1.eu}{jan@ju1.eu} \\
			\verb|\path{/home/foo/meindok.tex}| & \path{/home/foo/meindok.tex} \\
			\verb|\path{C:\TEMP\meindok.tex}| & \path{C:\TEMP\meindok.tex} \\
			\verb|\wort{Text} | & \wort{Text} \\
			\verb|\fremdwort{Text} |& \fremdwort{Text} \\
			\verb|Sonderzeichen: \& \% \$ \# \_ \{ \} |& \& \% \$ \# \_ \{ \} \\
			\verb|\LaTeX|& \LaTeX \\
			\verb|\dots|& \dots \\
		% -----------------------
		\bottomrule
	%\end{tabularx}
	\end{tabular}
	% =========================
		\caption{Logische Textauszeichnung}	% Caption anpassen!
		\label{tab:Textauszeichnung}	% Referenz anpassen!
	% =========================
\end{table}

\subsection{Punkte}\label{sec:Punkte}

(\autoref{tab:punkt} Punkte).
% Tabelle
\begin{table}[!hb] % hier einfügen
	\centering
	%\setlength{\tabcolsep}{5mm}      % Spaltenlänge fest
  \rowcolors{1}{}{lightgray!20}     % Farbe
	%\begin{tabularx}{\textwidth}{XX} % auto. Spaltenumbruch
	\begin{tabular} {ll}
		\toprule % Inhalt
		% -----------------------
			\verb|#| & \textbf{Beschreibung} \\
		\midrule
			\verb|Deutsch: Eins, zwei, ...| & Deutsch: Eins, zwei, ... \\
			\verb|Amerikanisch: One, two,~\dots| & Amerikanisch: One, two,~\dots \\
		% -----------------------
		\bottomrule
	%\end{tabularx}
	\end{tabular}
	% ===========================
		\caption{Punkte }	              % Caption anpassen!
		\label{tab:punkt}	              % Referenz anpassen!
	% ===========================
\end{table}

\subsection{Binde- und Gedankenstriche}\label{sec:Binde-Gedankenstriche}

(\autoref{tab:Bindestriche } Binde- und Gedankenstriche).
% Tabelle
\begin{table}[!hb]% hier
	\centering
	%\setlength{\tabcolsep}{5mm}      % Spaltenlänge fest
  \rowcolors{1}{}{lightgray!20}     % Farbe
	%\begin{tabularx}{\textwidth}{XX} % auto. Spaltenumbruch
	\begin{tabular} {ll}
		\toprule % Inhalt
		% -----------------------
			\verb|#| & \textbf{Beschreibung} \\
		\midrule
			\verb|n-zu-m-Abbildung| & n-zu-m-Abbildung \\
			\verb|11--19 Uhr| & 11--19 Uhr \\
			\verb|Berlin--Hamburg| & Berlin--Hamburg \\
			\verb|wahr -- oder falsch?| & wahr -- oder falsch? \\
			\verb|true---or false?| & true---or false? \\
			\verb|1, 0, $-$| & 1, 0, $-$ \\
		% -----------------------
		\bottomrule
	%\end{tabularx}
	\end{tabular}
	% ===========================
		\caption{Binde- und Gedankenstriche }% Caption anpassen!
		\label{tab:Bindestriche }	           % Referenz anpassen!
	% ===========================
\end{table}

\newpage % ++++++++++++++

\subsection{Listen}\label{sec:listen}

%% Liste
\begin{itemize}% Liste Punkt
	\item Text
	\item Text
\end{itemize}

\begin{enumerate}% Liste Aufzählung
	\item Text
	\item Text
\end{enumerate}

\begin{enumerate}% Liste 1.
	\item Text
	\begin{enumerate}% Liste a)
		\item Text
		\item Text
	\end{enumerate}% Liste 2.
	\item Text
	\begin{enumerate}% Liste b)
		\item Text
		\item Text
	\end{enumerate}
\end{enumerate}

\begin{enumerate}[label=(\roman*)]% (i), (ii), (iii) ...
	\item Text
	\item Text
\end{enumerate}

\begin{enumerate}[label={\arabic*\alph*)}]% 1a), 2b), 3c) ...
	\item Text
	\item Text
\end{enumerate}

\begin{enumerate}[label=\bfseries Punkt \Roman*]% Punkt I, Punkt II, Punkt III ...
	\item Text
	\item Text
\end{enumerate}

\section{Mathe - Beispiele}\label{mathe-bsp}

% \begin{align}	\quad \text{;} \quad \end{align}

\num{12345,678999}


\subsection{Potenzen }\label{sec:potenzen }\index{Potenzen }

allgemein:

\begin{align}
a^n &= a \cdot a \cdot ... \cdot a_n
\end{align}

Multiplikation: (gl.Basis, gl. Exponent)

\begin{align}
a^n \cdot a^m &= a^{n+m} \\
a^n \cdot b^n &= (a \cdot b)^n
\end{align}

Division:

\begin{align}
\frac{a^n}{a^m} &= a^{n-m} \\
a^{-n}          &= \frac{1}{a^n}   \\
a^0 						&= 1 \\
a^1				      &= a
\end{align}

Potenzen potenzieren:

\begin{align}
(a^n)^m         &= a^{n \cdot m} \\
\frac{a^n}{b^n} &= \left(\frac{a}{b}\right)^n
\end{align}

\begin{align}
a^b        &= e^{b \cdot ln \, a} \\
\sqrt[n]{a^m} &= a^{\frac{m}{n}}
\end{align}

(\autoref{code:mathe-out} Mathe).
% Quellcode
\lstset{language=[LaTeX]TeX} % C, [LaTeX]TeX, Bash, Python
\begin{lstlisting}[numbers=left, frame=l, framerule=0.1pt,%
% ======================
	caption={Quellcode in LaTeX, Mathe},% Anpassen!
	label={code:mathe-out}%
]%======================

	allgemein:

	\begin{align}
		a^n &= a \cdot a \cdot ... \cdot a_n
	\end{align}

	Multiplikation: (gl.Basis, gl. Exponent)

	\begin{align}
		a^n \cdot a^m &= a^{n+m} \\
		a^n \cdot b^n &= (a \cdot b)^n
	\end{align}

	Division:

	\begin{align}
		\frac{a^n}{a^m}   &= a^{n-m} \\
			a^{-n}          &= \frac{1}{a^n}   \\
			a^0 						&= 1 \\
			a^1				      &= a
	\end{align}

	Potenzen potenzieren:

	\begin{align}
		(a^n)^m         &= a^{n \cdot m} \\
		\frac{a^n}{b^n} &= \left(\frac{a}{b}\right)^n
	\end{align}

	\begin{align}
		a^b           &= e^{b \cdot ln \, a} \\
		\sqrt[n]{a^m} &= a^{\frac{m}{n}}
	\end{align}
\end{lstlisting}

\newpage % ++++++++++++++

\section{LaTeX - Befehle}

Textauszeichnung

(\autoref{code:textauszeichnung-out} Textauszeichnung).
% Quellcode
\lstset{language=[LaTeX]TeX} % C, [LaTeX]TeX, Bash, Python
\begin{lstlisting}[numbers=left, frame=l, framerule=0.1pt,%
% ======================
	caption={Quellcode in \LaTeX: Textauszeichnung},
	label={code:textauszeichnung-out}
]% =====================

	\emph{kursiv}
	\textrm{Antiqua}, \textsf{Grotesk}, \texttt{Maschinenschrift},
	\textmd{normal}, \textbf{breiter}, \textup{aufrecht}, \textsl{geneigt},
	\textit{kursiv}, \textsc{Kapitaelchen}
\end{lstlisting}

Schriftgroesse

(\autoref{code:schriftgroesse-out} Schriftgroesse).
% Quellcode
\lstset{language=[LaTeX]TeX} % C, [LaTeX]TeX, Bash, Python
\begin{lstlisting}[numbers=left, frame=l, framerule=0.1pt,%
% ======================
	caption={Quellcode in \LaTeX: Schriftgroesse},
	label={code:schriftgroesse-out}
]% =====================

  \tiny{winzig}, \scriptsize{sehr klein}, \footnotesize{klein},
	\small{klein}, \normalsize{normal}, \large{gross}, \Large{groesser},
	\LARGE{ganz gross}, \huge{riesig}, \Huge{gigantisch}
\end{lstlisting}

 eigene Befehle definieren

(\autoref{code:eigene-Befehle-definieren-out} eigene Befehle definieren).
% Quellcode
\lstset{language=[LaTeX]TeX} % C, [LaTeX]TeX, Bash, Python
\begin{lstlisting}[numbers=left, frame=l, framerule=0.1pt,%
% ======================
	caption={Quellcode in \LaTeX: eigene Befehle definieren},
	label={code:eigene-Befehle-definieren-out}
]% =====================

	\wort{Beispiel}
	\fremdwort{prezioes}
	\abstand{}

	\newcommand{\wort}[1]{\emph{#1}}
	\newcommand{\fremdwort}[1]{\textsf{#1}}
	\newcommand{\abstand}[1]{\vspace{5mm}{#1}}
	\newcommand{\wichtig}[2][red]{\textcolor{#1}{\emph{#2}}}

	quad, qquad, hspace{20mm}, vspace{20mm}
	Wichtig (Optionale Parameter)
	Wort Kursiv u. in Farbe
\end{lstlisting}

Eigene Umgebung

(\autoref{code:Eigene-Umgebung-out} Eigene Umgebung).
% Quellcode
\lstset{language=[LaTeX]TeX} % C, [LaTeX]TeX, Bash, Python
\begin{lstlisting}[numbers=left, frame=l, framerule=0.1pt,%
% ======================
	caption={Quellcode in \LaTeX: TEigene Umgebung},
	label={code:Eigene-Umgebung-out}
]% =====================

	Verwendung: \begin{hinweis}Ein Text.\end{hinweis}

	\newenvironment{hinweis}[1][Hinweis]{%
	  \begin{quote}
	  \color{meinblue}\rule{0.87\textwidth}{1pt}\\%
		\color{black}
	  \textbf{#1:}\\ %
	}{%
		\vspace{1mm}
	  \\ \color{meinblue}\rule[5ex]{0.87\textwidth}{1pt}%
	  \end{quote}
	}
\end{lstlisting}

farbige Infobox

(\autoref{code:farbigeInfobox-out} farbige Infobox).
% Quellcode
\lstset{language=[LaTeX]TeX} % C, [LaTeX]TeX, Bash, Python
\begin{lstlisting}[numbers=left, frame=l, framerule=0.1pt,%
% ======================
	caption={Quellcode in \LaTeX: farbige Infobox},
	label={code:farbigeInfobox-out}
]% =====================
	Anwendung: \myInfoBox{Text}

	\newcommand\myInfoBox[1]{%
	  \begin{quote}
		\fcolorbox{meinblue}{meingrey}{%
			\parbox{0.85\textwidth}{%
				\textbf{Hinweis:}\\%
				#1
			}
		}
	\end{quote}
	}
\end{lstlisting}




farbige Listenbox

(\autoref{code:farbigeListenbox-out} farbige Listenbox).
% Quellcode
\lstset{language=[LaTeX]TeX} % C, [LaTeX]TeX, Bash, Python
\begin{lstlisting}[numbers=left, frame=l, framerule=0.1pt,%
% ======================
	caption={Quellcode in \LaTeX: farbige Listenbox},
	label={code:farbigeListenbox-out}
]% =====================

	Anwendung:
	\myListenBox {
		\item Listenpunkt
		\item Listenpunkt
		\item Listenpunkt
	}

	\newcommand\myListenBox[1]{%
		\begin{quote}
			\fcolorbox{meinblue}{white}{%
				\parbox{0.85\textwidth}{%
					% Inhalt
					\textbf{Liste: }
					\begin{itemize}[label=$\square$]%checkbox
						#1
					\end{itemize}
				}
			}
		\end{quote}
	}
\end{lstlisting}

(\autoref{code:begriffe-out} Begriffe).
% Quellcode
\lstset{language=[LaTeX]TeX} % C, [LaTeX]TeX, Bash, Python
\begin{lstlisting}[numbers=left, frame=l, framerule=0.1pt,%
% ======================
	caption={Quellcode in \LaTeX: Begriffe},
	label={code:begriffe-out}
]% =====================

	Referenz: 				siehe~\ref{sec:abschnitt}.
	Zitat: 						siehe~\cite{Bos15}.
	Textauszeichnung: \wort{Beispiel}, \fremdwort{Fremdwort}
	Textabstand: 			\abstand{}
	Zahl Einheit: 		1\,l   z.\,B.
	nicht trennpaares Leerzeichen: ~
	Sonderzeichen: 		\& \% \$ \%
	Webadresse:				\url{http://www.LaTeXbuch.de}

	twoside=true		mschaltung Zweiseitig/Einseitiges Layout
	fontsize=12pt		Schriftgröße
	BCOR=10mm			  Bindekorrektur
	parskip					Gibt an wie neue Absätze gekennzeichnet werden.
								Hier empfohlene Beispielwerte:
									false:	Einzug der ersten Zeile
									half:		vertikaler Abstand von einer halben Zeile
									full:		vertikaler Abstand von einer Zeile
	paper=a4				DIN A4 Papier
	toc=listof			Im Inhaltsverzeichnis werden verzeichnisse wie Abbildungsverz.
								aufgenommen, wenn nicht gewünscht toc=nolistof
	toc=bib				  Literaturverzeichnis ohne Nummer im Inhaltsverzeichnis, oder nöchste zeile
	bibliography=totocnumbered	Literaturverzeichnis mit Nummer im Inhaltsverzeichnis, totoc ohne nummer
	open=right			Ein neues Kapitel fängt immer auf einer rechten Seite an, sonst open=any
	numbers=noenddot Nach DUDEN Werden Gliederungsnummern ohne Punkt am Ende gesetzt.
	headinclude			Kopfzeile zählt mit zum Graubereich der Seite
	headlines=2			zweizeilige Kopfzeile
	footexclude			Fußzeile enthält z.B. nur die Seitenzahl zählt deshalb nicht zum Graubereich der Seite
	pagesize=auto		Sorgt dafür, dass das PDF auch die richtige Größe hat
	version=last		welche Version des KOMA-Scripts verwendet werden soll
\end{lstlisting}

\newpage % ++++++++++++++

\section{Quellenangaben}

Quelle \textcite{schlosser_latex:2016}

Quelle Text\footfullcite{schlosser_latex:2016}

Quelle \cite{schlosser_latex:2016}

Raspberry Pi Kochbuch\footfullcite{monk_raspberry:2014} Quelle \cite{monk_raspberry:2014}

Elektronik-Hacks \footfullcite{monk_elektronik_hacks:2013} Quelle \cite{monk_elektronik_hacks:2013}

Action-Buch\footfullcite{monk_action_buch:2016} Quelle  \cite{monk_action_buch:2016}
% -----------------------
% -----------------------
% -----------------------

%\chapter{AdobeStockLizenz}
%%----------
% \section{ }
% \subsection{ }\label{ }\index{ }
%----------

%----
% \section{ }
% \subsection{ }\label{ }\index{ }
%----

\section{Einschränkungen der Standardlizenz von Adobe
Stock}\label{einschraenkungen-der-standardlizenz-von-adobe-stock}

Unter einer Standardlizenz zulässig:

\begin{itemize}% Liste Punkt% Liste Punkt
\item
  Reproduktion von bis zu 500.000 Kopien des Stockmediums auf
  Produktverpackungen und in gedruckten Marketingmaterialien sowie in
  digitalen Dokumenten oder Software.
\item
  Einschließen des Stockmediums in E-Mail-Marketing, mobile Werbung oder
  Fernsehprogramme, wenn weniger als 500.000 Zuschauer zu erwarten sind.
\item
  Veröffentlichen des Stockmediums auf einer Website ohne Einschränkung
  der Anzahl an Besucher. Wenn das Stockmedium unverändert auf einer
  Website der sozialen Medien veröffentlicht wird, ist eine
  Namensnennung erforderlich ((c) Autorenname -- stock.adobe.com).
\item
  Einschließen des Stockmediums in Produkte auf geringfügige Weise, z.
  B. in Lehrbüchern.
\end{itemize}

Unter einer Standardlizenz nicht zulässig:

Erstellen von Waren oder Handels- und Vertriebsprodukten, bei denen das
Stockmedium selbst den Hauptwert des Produkts ausmacht. So dürfen Sie
beispielsweise keine Poster, T-Shirts oder Kaffeetassen gestalten, die
hauptsächlich wegen des aufgedruckten Stockmediums gekauft werden.

Bilder

((C) Autorenname -- stock.adobe.com).

((C) Sergey Nivens -- stock.adobe.com).
%-----
%-----------

%\chapter{VorlageLaTeX}
%%----------
% \section{ }
% \subsection{ }\label{ }\index{ }
%----------

%----
% \section{ }
% \subsection{ }\label{ }\index{ }
%----

\section{Vorlage (LaTeX)}\label{vorlage-latex}

Modified source from v.14

\url{https://www.dcl.hpi.uni-potsdam.de/media/theses/}

\subsection{Literaturverzeichnis, Zitate und
Quellenangaben}\label{literaturverzeichnis-zitate-und-quellenangaben}

Tool \url{https://www.zotero.org/}

\subsection{Rechtschreibung \&
Grammatik}\label{rechtschreibung-grammatik}

\url{https://www.duden.de/Liste-der-rechtschreiblich-schwierigen-Woerter}

\subsection{Drucken \& Binden}\label{drucken-binden}

\url{http://www.potsdam-druck.de/digitaldruck/softcoverbuecher-drucken-binden.html}
%-----
%-----------

%
%\input{Projekt-files}
% --------------------------------------------------
%
% ggf. Anhang
%\appendix% ju - 25-Apr-2018  - anhang.tex 
% Modified source from
% <https://www.dcl.hpi.uni-potsdam.de/media/theses/>
%
\chapter{\appendixname}

\section{Eins}
Auch gibt es niemanden, der den Schmerz an sich liebt, sucht oder wünscht, nur, weil er Schmerz ist, es sei denn, es kommt zu zufälligen Umständen, in denen Mühen und Schmerz ihm große Freude bereiten können. Um ein triviales Beispiel zu nehmen, wer von uns unterzieht sich je anstrengender körperlicher Betätigung, außer um Vorteile daraus zu ziehen? Aber wer hat irgend ein Recht, einen Menschen zu tadeln, der die Entscheidung trifft, eine Freude zu genießen, die keine unangenehmen Folgen hat, oder einen, der Schmerz vermeidet, welcher keine daraus resultierende Freude nach sich zieht? Auch gibt es niemanden, der den Schmerz an sich liebt, sucht oder wünscht, nur, weil er Schmerz ist, es sei denn, es kommt zu zufälligen Umständen, in denen Mühen und Schmerz ihm große Freude bereiten können. Um ein triviales Beispiel zu nehmen, wer von uns unterzieht sich je anstrengender körperlicher Betätigung, außer um Vorteile daraus zu ziehen? Aber wer hat irgend ein Recht, einen Menschen zu tadeln, der die Entscheidung trifft, eine Freude zu genießen, die keine unangenehmen Folgen hat, oder einen, der Schmerz vermeidet, welcher keine daraus resultierende Freude nach sich zieht?Auch gibt es niemanden, der den Schmerz an sich liebt, sucht oder wünscht, nur,

\section{Zwei}
Auch gibt es niemanden, der den Schmerz an sich liebt, sucht oder wünscht, nur, weil er Schmerz ist, es sei denn, es kommt zu zufälligen Umständen, in denen Mühen und Schmerz ihm große Freude bereiten können. Um ein triviales Beispiel zu nehmen, wer von uns unterzieht sich je anstrengender körperlicher Betätigung, außer um Vorteile daraus zu ziehen? Aber wer hat irgend ein Recht, einen Menschen zu tadeln, der die Entscheidung trifft, eine Freude zu genießen, die keine unangenehmen Folgen hat, oder einen, der Schmerz vermeidet, welcher keine daraus resultierende Freude nach sich zieht? Auch gibt es niemanden, der den Schmerz an sich liebt, sucht oder wünscht, nur, weil er Schmerz ist, es sei denn, es kommt zu zufälligen Umständen, in denen Mühen und Schmerz ihm große Freude bereiten können. Um ein triviales Beispiel zu nehmen, wer von uns unterzieht sich je anstrengender körperlicher Betätigung, außer um Vorteile daraus zu ziehen? Aber wer hat irgend ein Recht, einen Menschen zu tadeln, der die Entscheidung trifft, eine Freude zu genießen, die keine unangenehmen Folgen hat, oder einen, der Schmerz vermeidet, welcher keine daraus resultierende Freude nach sich zieht?Auch gibt es niemanden, der den Schmerz an sich liebt, sucht oder wünscht, nur,

\section*{Drei (ohne extra Eintrag im Inhaltsverzeichnis)}
Auch gibt es niemanden, der den Schmerz an sich liebt, sucht oder wünscht, nur, weil er Schmerz ist,

\section*{Vier (ohne extra Eintrag im Inhaltsverzeichnis)}
Auch gibt es niemanden, der den Schmerz an sich liebt, sucht oder wünscht, nur, weil er Schmerz ist,
 % example

% Bibliographie
\ifisbook\cleardoubleemptypage\fi
\phantomsection\addcontentsline{toc}{chapter}{\refname}
\printbibliography[category=cited]
%\printbibliography[heading=bibintoc] % Bibliographie

% Eigenstaendigkeitserklaerung
%\ifisbook\pagestyle{plain}\cleardoubleemptypage% ju - 25-Apr-18 - erklaerung.tex 
% Modified source from
% <https://www.dcl.hpi.uni-potsdam.de/media/theses/>
%
\begin{otherlanguage}{ngerman}

\begin{center}\textsf{\textbf{Eidesstattliche Erklärung}}\end{center}
Hiermit versichere ich, dass meine Arbeit \enquote{\titel} selbständig verfasst wurde und dass keine anderen Quellen und Hilfsmittel als die angegebenen benutzt wurden. Diese Aussage trifft auch für alle Implementierungen und Dokumentationen im Rahmen dieses Projektes zu.\\

\noindent
\ort, den \datum,
\vspace{2cm}

\begin{center}
\begin{tabular}{C{6cm}}
\rowcolor{white}
\hline
{\small({\autor})}
\end{tabular}
\end{center}

\end{otherlanguage}
\fi % example

\end{document}
%============================
%============================
%============================
